\documentclass[11pt,a4paper]{article}
\usepackage[T1]{fontenc}
\usepackage[utf8]{inputenc}
\usepackage{lmodern}
\usepackage[margin=25.5mm,headheight=15pt]{geometry}
\usepackage{amsmath,amssymb,amsthm,mathtools}
\usepackage{booktabs,longtable,tabularx,array,microtype,enumitem,xcolor}
\usepackage{fancyhdr,listings,xurl,titlesec}
\definecolor{navy}{RGB}{30,57,82}
\definecolor{ink}{RGB}{32,38,44}
\definecolor{pale}{RGB}{239,243,247}
\usepackage[colorlinks=true,linkcolor=ink,citecolor=ink,urlcolor=navy,
 bookmarksnumbered=true]{hyperref}
\usepackage{bookmark}
\hypersetup{pdftitle={Coarse Degrees with Unattained Turing Infima},
 pdfauthor={Research note prepared for Vladimir Reshetnikov},
 pdfsubject={Explicit counterexample, perfect minimal pairs, and least Turing representatives}}
\urlstyle{same}
\titleformat{\section}{\large\bfseries\color{navy}}{\thesection}{0.7em}{}
\titleformat{\subsection}{\normalsize\bfseries\color{navy}}{\thesubsection}{0.65em}{}
\setlength{\parskip}{3.5pt}
\setlength{\parindent}{1.2em}
\setlength{\emergencystretch}{3em}
\setlist{nosep,leftmargin=*}
\pagestyle{fancy}\fancyhf{}
\fancyhead[L]{\small\color{navy}Coarse degrees and Turing infima}
\fancyhead[R]{\small September 2026}
\fancyfoot[C]{\thepage}
\renewcommand{\headrulewidth}{0.3pt}
\newcommand{\N}{\mathbb N}
\newcommand{\Cantor}{2^{\N}}
\newcommand{\Baire}{\N^{\N}}
\newcommand{\leT}{\leq_{\mathrm T}}
\newcommand{\geT}{\geq_{\mathrm T}}
\newcommand{\eqT}{\equiv_{\mathrm T}}
\newcommand{\degT}{\operatorname{deg}_{\mathrm T}}
\newcommand{\CD}{\operatorname{CD}}
\newcommand{\nc}{\mathrm{nc}}
\newcommand{\uc}{\mathrm{uc}}
\newcommand{\Spec}{\mathcal S}
\newcommand{\Core}{\mathcal K}
\newcommand{\Res}{\mathcal R}
\newcommand{\Block}{\mathcal B}
\newcommand{\zerodeg}{\mathbf 0}
\newcommand{\zdeg}{\mathbf z}
\newcommand{\udens}{\overline\rho}
\newcommand{\concat}{\mathbin{{}^\frown}}
\newcommand{\dom}{\operatorname{dom}}
\newcommand{\Graph}{\operatorname{Graph}}
\newtheorem{theorem}{Theorem}[section]
\newtheorem{lemma}[theorem]{Lemma}
\newtheorem{proposition}[theorem]{Proposition}
\newtheorem{corollary}[theorem]{Corollary}
\theoremstyle{definition}
\newtheorem{definition}[theorem]{Definition}
\theoremstyle{remark}
\newtheorem{remark}[theorem]{Remark}
\lstset{basicstyle=\small\ttfamily,columns=fullflexible,breaklines=true,
 frame=single,rulecolor=\color{navy},backgroundcolor=\color{pale},
 showstringspaces=false,xleftmargin=0pt,aboveskip=9pt,belowskip=9pt}
\pdfstringdefDisableCommands{\def\leT{<=T}\def\N{N}\def\zerodeg{0}}
\begin{document}
\hypersetup{pageanchor=false}
\begin{titlepage}
\vspace*{12mm}
{\small\sffamily\bfseries\color{navy}COMPUTABILITY THEORY\quad /\quad RESEARCH NOTE\par}
\vspace{15mm}
{\Huge\bfseries\color{navy}Coarse Degrees with\par Unattained Turing Infima\par}
\vspace{8mm}
{\Large An explicit c.e. counterexample,\par perfect minimal pairs, and relativization\par}
\vspace{13mm}
\noindent\rule{\textwidth}{0.5pt}
\vspace{5mm}
\noindent Prepared for Vladimir Reshetnikov\\[3pt]
18 September 2026
\vspace{10mm}
\begin{abstract}
We investigate whether a nonuniform coarse-equivalence class must contain a
representative of least Turing degree. We give an explicit computably enumerable
set whose coarse descriptions contain a perfect family of pairwise Turing
minimal pairs. Its uniform and nonuniform coarse classes consequently have no
least Turing representative, even when arbitrary total numerical functions are
allowed as representatives. Two witnessing descriptions can be chosen below
$\emptyset''$. The construction uses individually computable protected columns
and a two-reservoir finite-extension lemma; all infinite arguments are proved
in English. Relativizing and adding robust block coding yields, for every Turing
degree $\zdeg$, a coarse class whose representative spectrum has greatest lower
bound $\zdeg$ but does not attain it. We also characterize the nonuniform classes
having a least representative as precisely the classes of robustly coded reals.
The bare negative answer is already implicit in a 2015 result of
Hirschfeldt--Jockusch--Kuyper--Schupp; it is not claimed here as a new discovery.
The originality of the stronger perfect-family and prescribed-infimum
formulations remains unverified.
\end{abstract}
\vfill
\noindent\small Conventional proofs, a concrete enumeration program, and finite checks
are supplied. No Lean formalization or independent referee verification is claimed.
\end{titlepage}
\hypersetup{pageanchor=true}
\setcounter{tocdepth}{1}
\tableofcontents
\vspace{8mm}
\noindent\colorbox{pale}{\parbox{\dimexpr\textwidth-2\fboxsep\relax}{
\textbf{Reading route.} Section~\ref{sec:results} states the result and its
provenance. Sections~\ref{sec:columns}--\ref{sec:mainproof} give the independent
counterexample proof. Section~\ref{sec:relative} proves the prescribed-infimum
extension. Section~\ref{sec:classification} explains exactly where least
representatives do occur.}}
\newpage

\section{The selected question and the result}\label{sec:results}
The target is entry C1 of the supplied research report~\cite{Survey}:
\begin{equation}\label{eq:question}
 (\forall f\in\Baire)(\exists g\equiv_{\nc}f)
 (\forall h\equiv_{\nc}f)\;g\leT h.
\end{equation}
It is the nonuniform coarse instance of Gerdes's Question~7~\cite{Gerdes}.
The same least-representative question is displayed in the author's July 2026
ASL slides~\cite{GerdesSlides}. Thus the choice is based on a specifically
stated question, not a conjecture invented to fit an easy construction.

Here $X\approx Y$ means that their disagreement set has asymptotic density
zero. The identity functional makes any two such functions uniformly coarsely
equivalent. The decisive distinction is between the Turing degrees of
representatives and their common coarse degree.

\begin{theorem}[Explicit counterexample]\label{thm:main}
There is a computably enumerable set $A\subseteq\N$ with no computable coarse
description and a nonempty perfect set $\mathcal P\subseteq\Cantor$ such that:
\begin{enumerate}[label=\textup{(\roman*)}]
\item every $X\in\mathcal P$ satisfies $X\approx A$;
\item if $X,Y\in\mathcal P$ are distinct and a total function $F$ satisfies
$F\leT X$ and $F\leT Y$, then $F$ is computable;
\item there are distinct $X_0,X_1\in\mathcal P$ with $X_0,X_1\leT\emptyset''$.
\end{enumerate}
Consequently both the uniform and the nonuniform coarse class of $A$ have no
least Turing representative. This remains true with representatives in
$\Baire$, and remains true after restricting their Turing degrees to be at
most $\degT(\emptyset'')$.
\end{theorem}

Every member of $\mathcal P$ is noncomputable. Thus (ii) says that every two
distinct members form a minimal pair in the \emph{ordinary Turing degrees}.
They are not a minimal pair of coarse degrees: their coarse degrees are equal
and nonzero. In particular, the example is stronger than merely having two
Turing-incomparable representatives.

The next theorem identifies an additional feature of the representative
spectrum. The infimum is taken in the full partial order of Turing degrees,
not in a linear ordering of complexity classes.

\begin{theorem}[Any prescribed degree can be an unattained infimum]
\label{thm:infimum}
For every real $Z$ there is a binary set $A_Z$, c.e. in $Z$, such that for
$r\in\{\nc,\uc\}$ the spectrum
\[
 \Spec_r(A_Z)=\{\degT(F):F\in\Baire,\ F\equiv_r A_Z\}
\]
has greatest lower bound $\degT(Z)$, but does not contain $\degT(Z)$.
Moreover, $\CD(A_Z)$ contains a perfect family whose distinct members have
Turing-degree meet exactly $\degT(Z)$. Two such members can be chosen below
$Z''$.
\end{theorem}

\subsection{What is, and is not, new here}
Hirschfeldt, Jockusch, Kuyper, and Schupp define $X^{\mathfrak c}$ to be the
sets computable from every coarse description of $X$. Their Theorem~4.2
states that $X^{\mathfrak c}=\zerodeg$ for a 1-generic $X$, which is not
coarsely computable. Their Theorem~4.3 gives the same conclusion when the
coarse computability bound $\gamma(X)$ is $1$~\cite{HJKS}.

The implication for~\eqref{eq:question} is immediate but worth spelling out.
If $g$ were a least representative of the coarse class of such a 1-generic
$X$, then $g$ would be computable from every coarse description of $X$,
since every one of those descriptions belongs to the same coarse class.
Applying the stated theorem to the graph of $g$ would make $g$ computable.
The reduction $X\le_{\nc}g$ would then give a computable coarse description
of $X$, a contradiction. The same argument works for uniform coarse classes.

Accordingly, the bare negative answer is an implication of established
literature, rather than a priority claim of this note. What is developed
below is a self-contained construction of stronger witnesses and the
prescribed-infimum extension. These formulations were not located in the
targeted source search, but that does not establish their originality.
The proof does not use genericity, randomness, or the cited cone-avoiding
compactness theorem as an unproved dependency.

\section{Definitions and the least-degree obstruction}\label{sec:definitions}
For $S\subseteq\N$ and $N>0$, put
\[
 \rho_N(S)=\frac{|S\cap[0,N)|}{N},\qquad
 \udens(S)=\limsup_{N\to\infty}\rho_N(S).
\]
We identify a set with its total characteristic function. For numerical
functions $f,g$, write $f\triangle g=\{n:f(n)\ne g(n)\}$, and put
\[
 f\approx g\ \Longleftrightarrow\ \udens(f\triangle g)=0,
 \qquad \CD(f)=\{D\in\Baire:D\approx f\}.
\]
Thus descriptions are total; they may give incorrect answers, but only on
a density-zero set. These are the coarse-description conventions of
Gerdes~\cite{Gerdes}. We allow numerical descriptions even when the target
is binary.

Define
\begin{align*}
 f\le_{\nc}g
 &\Longleftrightarrow
 (\forall D\in\CD(g))(\exists E\in\CD(f))\;E\leT D,\\
 f\le_{\uc}g
 &\Longleftrightarrow
 (\exists\Phi)(\forall D\in\CD(g))\;\Phi^D\in\CD(f).
\end{align*}
In the second line the functional must be total on each valid input oracle.
Equivalence is mutual reducibility. The coarse degree containing the computable functions is
called the zero coarse degree; its representatives are exactly the
coarsely computable functions.

\begin{lemma}[Descriptions stay in the class]\label{lem:descriptions}
If $D\approx f$, then $\CD(D)=\CD(f)$ and $D\equiv_{\uc}f$.
If $f$ has no computable coarse description, then neither its uniform nor
its nonuniform coarse class contains a computable function.
\end{lemma}
\begin{proof}
The union of two density-zero sets has density zero: its counting function
is at most the sum of the two counting functions. Since
$E\triangle f\subseteq(E\triangle D)\cup(D\triangle f)$, the relation
$\approx$ is transitive. It is also symmetric. This proves equality of the
description classes. The identity functional supplies both uniform reductions.

Suppose a computable $g$ is coarsely equivalent to $f$. It is a coarse
description of itself. Applying $f\le_{\nc}g$ to that input gives a member
of $\CD(f)$ computable from $g$, hence computable, a contradiction.
A uniform reduction is in particular a nonuniform one.
\end{proof}

\begin{lemma}[Minimal-pair certificate]\label{lem:certificate}
Suppose $A$ has no computable coarse description and $X,Y\in\CD(A)$
have only computable total functions as common Turing lower bounds.
Then $[A]_{\nc}$ and $[A]_{\uc}$ have no least Turing representative.
\end{lemma}
\begin{proof}
Both $X$ and $Y$ are representatives by Lemma~\ref{lem:descriptions}.
A least representative $g$ would satisfy $g\leT X,Y$, and hence would be
computable. Lemma~\ref{lem:descriptions} excludes such a representative.
Nothing in this argument requires $g$ to be binary.
\end{proof}

\begin{remark}[Binaryization]
For a numerical function $D$, set $\beta(D)(n)=1$ when $D(n)=1$, and set
it to $0$ otherwise. If $A$ is binary, then
$\beta(D)\triangle A\subseteq D\triangle A$.
Consequently a computable numerical coarse description of a binary set
would give a computable binary one. The same observation holds relative to
any fixed oracle.
\end{remark}

\section{The explicit c.e. diagonal set}\label{sec:columns}
Fix an effective enumeration $(\varphi_e)_{e\in\N}$ of all partial
computable numerical functions. No totality test is built into this
enumeration. Define
\begin{equation}\label{eq:columns}
 p_e(j)=2^e(2j+1)-1,\qquad R_e=p_e[\N].
\end{equation}
The sets $R_e$ partition $\N$: $e$ is the exponent of $2$ in $n+1$.
The first three columns are the even numbers, the numbers congruent to
$1$ modulo $4$, and the numbers congruent to $3$ modulo $8$.
Put
\begin{equation}\label{eq:protected}
 P_k=\bigcup_{e<k}R_e,\qquad Q_k=\N\setminus P_k
       =\{n:2^k\mid n+1\}.
\end{equation}
For every $N>0$,
\begin{equation}\label{eq:tailcount}
 |Q_k\cap[0,N)|=\left\lfloor\frac{N}{2^k}\right\rfloor,
 \qquad \rho(R_e)=2^{-(e+1)}.
\end{equation}
In particular every $Q_k$ is infinite, while $\rho(Q_k)=2^{-k}$ tends to zero.

Define $A$ by the following positive convergence condition:
\begin{equation}\label{eq:A}
\begin{split}
 A(p_e(j))=1\quad\Longleftrightarrow\quad &
 \varphi_e(p_e(j))\downarrow=0\ \text{ and }\\[-2pt]
 & (\forall i\le j)\;
       \varphi_e(p_e(i))\downarrow\in\{0,1\}.
\end{split}
\end{equation}
At all other positions $A$ has value zero. This is a definition of one
particular set after choosing the enumeration, not a stage-dependent
choice of a future witness.

\begin{lemma}[Computable finite unions of columns]\label{lem:columncomp}
The set $A$ is c.e. For every fixed $e$, $A\cap R_e$ is computable, and
for every fixed $k$, $A\cap P_k$ is computable. A computable uniform
sequence of indices for these sets is not asserted.
\end{lemma}
\begin{proof}
Membership in $A$ has a finite positive witness: finitely many computations
must have halted with specified allowed values. Dovetail these computations
and enumerate $p_e(j)$ once the condition in~\eqref{eq:A} has been witnessed.
This enumerates exactly $A$ and never retracts an enumerated point.

Fix $e$. There are two cases. If every $\varphi_e(p_e(i))$ converges with a
binary value, then on $R_e$ the set $A$ is the complement of this total
computable sequence. To decide membership, first test whether the input
belongs to $R_e$; if so, run the corresponding computation and complement
its value. The computation terminates in this case.

Otherwise there is a least $q$ at which the computation either diverges
or converges with a nonbinary value. For every $j\ge q$, the second
conjunct in~\eqref{eq:A} fails. Thus $A\cap R_e$ is contained in the finite
set $\{p_e(0),\ldots,p_e(q-1)\}$, and is computable as a finite set.
These two cases prove individual computability without deciding which
case holds effectively from $e$. A finite union of the first $k$ such
computable sets gives $A\cap P_k$.
\end{proof}

\begin{lemma}[Diagonalization and close computable approximants]
\label{lem:noncoarse}
The set $A$ has no computable coarse description. Nevertheless, for every
$k$ there is a computable set $C_k$ with
$\udens(A\triangle C_k)\le2^{-k}$.
\end{lemma}
\begin{proof}
Let $C$ be any computable binary set. Choose an index $e$ for its total
characteristic function. Every computation $\varphi_e(p_e(i))$ then
converges with a binary value. Formula~\eqref{eq:A} makes
$A(p_e(j))=1-C(p_e(j))$ for every $j$. Consequently
\[
 R_e\subseteq A\triangle C,
 \qquad
 \liminf_N\rho_N(A\triangle C)\ge 2^{-(e+1)}>0.
\]
No computable binary set is a coarse description of $A$. Binaryization
excludes computable numerical descriptions as well.

For the second assertion take $C_k=A\cap P_k$, which is computable by
Lemma~\ref{lem:columncomp}. Its disagreement with $A$ is contained in
$Q_k$, so~\eqref{eq:tailcount} gives the required bound.
\end{proof}

\paragraph{The useful tension.}
Each finite collection of positive-density columns is computable, yet the
whole set defeats every computable coarse description. The finite-extension
construction will use one finite collection at a time. It must not be given
a fictitious uniform ordinary procedure for computing all the protected
column values simultaneously. The next lemma makes this distinction explicit.

\section{A two-reservoir computation lemma}\label{sec:reservoir}
We formulate the construction relative to a fixed oracle $Z$ so that it
can be used again later. The unrelativized application has $Z=\emptyset$.
Let $A\subseteq\N$, let $P_k$ be increasing $Z$-computable sets, and
suppose that $A\cap P_k$ is $Z$-computable for each fixed $k$.
For a finite binary string $\sigma$, define
\begin{equation}\label{eq:reservoir}
 \Res(\sigma,k)=\{X\in\Cantor:X\succeq\sigma\ \text{ and }\
  (\forall n\ge|\sigma|)\,[n\in P_k\Rightarrow X(n)=A(n)]\}.
\end{equation}
Such a reservoir is nonempty: extend $\sigma$ with the prescribed bits on
$P_k$ and, for example, zeros at the other positions.

An extension $\alpha\succeq\sigma$ is \emph{admissible at level $k$} if
it agrees with $A$ on $P_k\cap[|\sigma|,|\alpha|)$. With the fixed
indices for $P_k$ and $A\cap P_k$, this is a $Z$-decidable test.
It does not query the full noncomputable set $A$.

If $\alpha$ is admissible, then
\begin{equation}\label{eq:shrink}
 \Res(\alpha,k)\subseteq\Res(\sigma,k),\qquad
 \Res(\alpha,k+1)\subseteq\Res(\sigma,k).
\end{equation}
For the second inclusion, the old protected values between the two prefix
lengths were already respected, and beyond the longer prefix they are
still protected because $P_k\subseteq P_{k+1}$. Values fixed in the prefix
are never retroactively changed.

We use finite-oracle computations $\Phi_e^{Z\oplus\alpha}(n)\downarrow=v$.
This means that the computation halts with value $v$ and all queries to the
second oracle have positions below $|\alpha|$. The first oracle $Z$ remains
fully available. Such convergence is a $Z$-c.e. condition and persists
under extension of $\alpha$.

\begin{lemma}[Cross-disagreement or computable common output]
\label{lem:product}
Fix $k$, two prefixes $\sigma,\tau$, and two functional indices $e,j$.
Exactly one of the following alternatives holds.
\begin{enumerate}[label=\textup{(\arabic*)}]
\item There are admissible extensions $\alpha\succeq\sigma$ and
$\beta\succeq\tau$ and an input $n$ for which
\[
 \Phi_e^{Z\oplus\alpha}(n)\downarrow
 \ne \Phi_j^{Z\oplus\beta}(n)\downarrow.
\]
The disagreement holds for every pair of full oracles in the two smaller
reservoirs.
\item No such extensions exist. Whenever $X\in\Res(\sigma,k)$ and
$Y\in\Res(\tau,k)$ yield a common total function
$F=\Phi_e^{Z\oplus X}=\Phi_j^{Z\oplus Y}$, one has $F\leT Z$.
\end{enumerate}
Existence in alternative \textup{(1)} is a $\Sigma^0_1(Z)$ condition
given the fixed reservoir indices.
\end{lemma}
\begin{proof}
If the extensions in (1) exist, finite use preserves the displayed unequal
answers for all further extensions. There can be no common total function
for those smaller reservoirs.

Suppose no cross-disagreement exists, and suppose the $X,Y,F$ in (2)
are given. We describe a $Z$-oracle computation of $F(n)$. Enumerate all
admissible finite extensions $\alpha$ of $\sigma$, and dovetail the
computations $\Phi_e^{Z\oplus\alpha}(n)$. Output the value of the first
one that halts. The search terminates: a sufficiently long initial segment
of the actual $X$ is admissible and witnesses its convergent computation
on $n$.

To verify correctness, take a sufficiently long initial segment $\beta$
of the actual $Y$ witnessing $\Phi_j^{Z\oplus\beta}(n)=F(n)$. It is
admissible. Any answer different from $F(n)$ found by the search would
constitute the forbidden cross-disagreement with this $\beta$. Hence the
search returns exactly $F(n)$.

The procedure uses only $Z$, fixed finite prefixes, fixed indices for the
$Z$-computable protected sets, and the functional index $e$. These finitely
many integers may be hardcoded. Neither $X,Y$ nor the full set $A$ is an
oracle for the procedure. Finally, the existence assertion in (1) has
finite witnesses and a $Z$-c.e. verification, proving its complexity bound.
\end{proof}

\begin{remark}[The uniformity boundary]
The lemma proves that a common \emph{total} output is $Z$-computable.
It does not claim that the output exists, that either functional is total
on a reservoir, or that a single ordinary algorithm can decide which
alternative applies. For $Z=\emptyset$, the latter decision may use
$\emptyset'$. Its role in selecting the construction is different from
the ordinary computability of a common output in alternative (2).
\end{remark}

\section{Perfect fusion with vanishing error}\label{sec:fusion}
The preceding lemma yields a useful general theorem. Its assumptions are
stronger than merely being arbitrarily close in density to computable sets:
they provide nested regions on which the target is exactly computable.

\begin{theorem}[Computable-reservoir fusion]\label{thm:fusion}
Fix $Z$ and $A\subseteq\N$. Suppose $(P_k)$ is an increasing sequence of
$Z$-computable sets, each $A\cap P_k$ is $Z$-computable, and
\[
 Q_k=\N\setminus P_k\text{ is infinite for every }k,
 \qquad \lim_{k\to\infty}\udens(Q_k)=0.
\]
No uniform sequence of indices is required. There is a perfect set
$\mathcal P\subseteq\CD(A)$ such that, whenever $X,Y\in\mathcal P$ are
distinct and $F$ is total,
\begin{equation}\label{eq:relativepair}
 F\leT Z\oplus X\ \text{ and }\ F\leT Z\oplus Y
       \quad\Longrightarrow\quad F\leT Z.
\end{equation}
There are lengths $L_k$ such that simultaneously for all $X\in\mathcal P$,
\begin{equation}\label{eq:uniformerror}
 X\triangle A\subseteq[0,L_k)\cup Q_k.
\end{equation}
If an oracle $W\geT Z'$ computes indices for $P_k$ and $A\cap P_k$
uniformly in $k$, the finite level construction is $W$-computable.
A member indexed by a branch $U$ is then computable from $W\oplus U$.
\end{theorem}
\begin{proof}
At level $s$ construct strings $\sigma_\eta$ for $\eta\in2^s$, all of one
length $L_s$, with distinct strings incompatible. They represent the
reservoirs $\Res(\sigma_\eta,s)$. Start with the empty string and $L_0=0$.

\emph{Splitting.} From each level-$s$ node choose a fresh position
$q\in Q_s$ beyond $L_s$. Such a position exists because $Q_s$ is infinite.
Extend admissibly through this position in two ways, giving $q$ the values
$0$ and $1$. Use the prescribed values on $P_s$ at the intervening
positions, and zeros elsewhere. This creates children labelled
$\eta\concat0$ and $\eta\concat1$. Children of different parents remain
incompatible because their parents were incompatible.

\emph{Meeting the finite collection of pair requirements.}
Consider every ordered pair of distinct labels
$\xi,\zeta\in2^{s+1}$ and every $e,j\le s$. There are finitely many such
tasks, so fix any finite ordering of them. For the current prefixes at
$\xi,\zeta$, apply Lemma~\ref{lem:product} at protection level $s$.
If cross-disagreement exists, replace both prefixes by witnesses to it.
If not, leave the prefixes unchanged and record alternative (2).
Other prefixes are unchanged in this task.

Every extension made in this process is admissible. Previously established
disagreement is permanent. A previous no-disagreement conclusion is also
preserved, since the reservoirs only shrink. Thus a later task cannot
injure an earlier one. Each level requires only finitely many decisions
and finite extensions; there is no infinite process inside a single level.

\emph{Raising protection.} Once the tasks are complete, pad all prefixes
admissibly to a common length $L_{s+1}>L_s$, with $L_{s+1}\ge s+1$.
Now use protection level $s+1$ beyond these prefixes. By~\eqref{eq:shrink},
every child reservoir is contained in its parent reservoir and in the
reservoirs used by its completed tasks. Raising protection never invalidates
a prefix, because it prescribes only positions beyond that prefix.
This finishes the induction.

\emph{The perfect family.} For $U\in\Cantor$ put
\[
 X_U=\bigcup_s\sigma_{U\upharpoonright s}.
\]
The union defines a total binary sequence since $L_s\to\infty$.
Different $U$'s give different sequences, by the splitting step.
Let $\mathcal P=\{X_U:U\in\Cantor\}$. Equivalently,
\[
 \mathcal P=\bigcap_s\ \bigcup_{\eta\in2^s}[\sigma_\eta],
\]
where $[\sigma]$ is the cylinder of extensions of $\sigma$.
The right-hand side is closed. Each of its points determines unique,
compatible labels at all levels and therefore is one of the $X_U$'s.
Given a member and any finite initial segment of it, take a sufficiently
late level whose length covers that segment and then split at a later
level. Another branch shares the segment but gives a different member.
Thus the closed set has no isolated points, and is perfect.

\emph{Density convergence.} Fix $k$. All later reservoirs are contained
in the level-$k$ reservoir of the relevant branch. Consequently no member
can disagree with $A$ on $P_k$ beyond position $L_k$.
This proves~\eqref{eq:uniformerror}. For every $N>0$,
\begin{equation}\label{eq:densitybudget}
 \rho_N(X_U\triangle A)\le \frac{L_k}{N}+\rho_N(Q_k).
\end{equation}
First let $N\to\infty$, with $k$ fixed, to obtain
$\udens(X_U\triangle A)\le\udens(Q_k)$. Then let $k\to\infty$.
The upper density is zero. In particular every $X_U$ lies in $\CD(A)$.
The order of these two limiting operations is essential.

\emph{Common lower functions.} Suppose $U\ne V$ and
$F=\Phi_e^{Z\oplus X_U}=\Phi_j^{Z\oplus X_V}$ is total.
Choose $s$ sufficiently large that $e,j\le s$ and the two branches have
distinct labels at level $s+1$. Their ordered-pair task was processed at
stage $s$. The disagreement alternative would give different answers at
one input for the actual final oracles, contrary to equality with $F$.
The task therefore used the no-disagreement alternative. The final oracles
still belong to its reservoirs, so Lemma~\ref{lem:product} gives $F\leT Z$.
Every pair of functional indices is covered this way, proving
\eqref{eq:relativepair}.

\emph{Effectivity.} Given the stated index sequence, $W$ can carry out
admissibility tests, finite padding, and searches for fresh free positions.
The cross-disagreement question is $\Sigma^0_1(Z)$, and hence decidable
using $Z'\leT W$. When the answer is positive, enumerate witnesses and
choose the first in a fixed effective ordering. These rules determine
the entire finite level system using $W$. With a branch oracle $U$, find
a level of length greater than the requested position and read its bit.
Thus $X_U\leT W\oplus U$.
\end{proof}

\paragraph{Why ordinary Cohen forcing alone would not suffice.}
Unrestricted finite extensions can diagonalize computations but can also
introduce errors of positive density. Permanently fixing all of $A$ outside
a single computable density-zero set would leave too much noncomputable
information in each reservoir. Here the protected part is computable at
each stage, while the density of the unprotected tail tends to zero only
across stages. This is what lets the computational and density requirements
coexist.

\section{Completing the counterexample and bounding its witnesses}
\label{sec:mainproof}
Apply Theorem~\ref{thm:fusion} with $Z=\emptyset$, the set $A$ of
\eqref{eq:A}, and the sets $P_k$ of~\eqref{eq:protected}.
Lemma~\ref{lem:columncomp} supplies the protected computations, and
\eqref{eq:tailcount} supplies the tail hypothesis. In this instance the
error estimate is the explicit uniform bound
\begin{equation}\label{eq:explicitbudget}
 \rho_N(X\triangle A)
 \le\frac{L_k}{N}+\frac{\lfloor N/2^k\rfloor}{N}
 \le\frac{L_k}{N}+2^{-k}
 \qquad(X\in\mathcal P).
\end{equation}
The same containment and estimate hold for $X\triangle Y$ whenever
$X,Y\in\mathcal P$, since beyond $L_k$ both agree with $A$ on $P_k$.

No member of $\mathcal P$ is computable, by Lemma~\ref{lem:noncoarse}.
The common-lower-function conclusion therefore supplies a genuine minimal
pair from any two distinct members. In particular the family represents
continuum many pairwise incomparable Turing degrees in a single coarse class.

\begin{lemma}[An index selector below the second jump]\label{lem:selector}
For the set $A$ of~\eqref{eq:A}, a sequence of ordinary computable indices
for $A\cap P_k$ can be obtained uniformly using $\emptyset''$.
\end{lemma}
\begin{proof}
For each $e$, the assertion
\[
 (\forall i)(\exists s)\;
 [\varphi_e(p_e(i))\text{ halts by stage }s\text{ with a binary value}]
\]
is $\Pi^0_2$ and is decidable by $\emptyset''$.
If it holds, a program for $A\cap R_e$ tests membership in $R_e$ and
complements the indicated convergent computation, as in
Lemma~\ref{lem:columncomp}. This is an ordinary program with $e$ as a
parameter, not a program using the second jump.

If the assertion fails, use $\emptyset'$ to find the first index $q$ at
which convergence with a binary value fails. Such failure at a fixed
index is decidable by $\emptyset'$. Run the finitely many successful
computations below $q$, and output an index for the resulting finite set.
The second jump can perform all these steps. Finally combine the first
$k$ indices by finite union. The resulting index is again for an ordinary
computable set; the oracle was used to choose that index.
\end{proof}

\begin{proof}[Proof of Theorem~\ref{thm:main}]
The preceding application of the fusion theorem proves (i) and (ii).
Lemma~\ref{lem:selector} and the effective clause of
Theorem~\ref{thm:fusion} make the finite level system
$\emptyset''$-computable. Use, for example, the two computable branch
labels $0^\omega$ and $1^\omega$. They give distinct members
$X_0,X_1\leT\emptyset''$, proving (iii).
Lemma~\ref{lem:certificate} rules out a least representative in either
full coarse class.

For the restricted statement, the same two representatives still belong
to the spectrum below $\degT(\emptyset'')$. A least representative of
that restricted spectrum would be below both, hence computable, and the
same contradiction applies. This is not an assertion that \emph{every}
member of the perfect family is below $\emptyset''$; an arbitrary branch
label may carry additional information.
\end{proof}

\section{Prescribing the unattained infimum}\label{sec:relative}
We now keep an arbitrary oracle $Z$ fixed. Two ingredients are needed:
a relative version of the diagonal set and a coding that every coarse
description can decode.

\subsection{Relativized diagonal columns}
Use an effective enumeration $(\varphi_e^Z)$ of the partial $Z$-computable
functions and define $B_Z$ by
\begin{equation}\label{eq:BZ}
\begin{split}
 B_Z(p_e(j))=1\quad\Longleftrightarrow\quad &
 \varphi_e^Z(p_e(j))\downarrow=0\ \text{ and }\\[-2pt]
 &(\forall i\le j)\;
       \varphi_e^Z(p_e(i))\downarrow\in\{0,1\}.
\end{split}
\end{equation}
Every step of Lemmas~\ref{lem:columncomp} and~\ref{lem:noncoarse}
relativizes: $B_Z$ is c.e. in $Z$, each $B_Z\cap P_k$ is
$Z$-computable, and $B_Z$ has no $Z$-computable coarse description.
For completeness, a total $Z$-computable candidate has an index $e$;
on the entire positive-density column $R_e$, formula~\eqref{eq:BZ}
complements it. The other column case is still a finite set.
The proof of Lemma~\ref{lem:selector} likewise gives a $Z''$-computable
sequence of indices for the protected $Z$-computable pieces.

\subsection{A robust block code}
The following elementary code repeats each bit on a block occupying a
fixed positive proportion of its ending prefix. It is a standard type of
coarsening; we give its full proof and make no novelty claim for this ingredient.
Set
\begin{equation}\label{eq:block}
 b_n=2^{2^n},\qquad J_n=[b_n,b_{n+1}),\qquad
 \Block(Z)(m)=
 \begin{cases}
 0,&m<2,\\
 Z(n),&m\in J_n.
 \end{cases}
\end{equation}

\begin{lemma}[Every coarse description decodes the block code]
\label{lem:block}
One has $\Block(Z)\eqT Z$. Every $D\in\CD(\Block(Z))$ computes $Z$.
The latter assertion is nonuniform in the description $D$.
\end{lemma}
\begin{proof}
A $Z$-oracle computes the block containing any input and hence its bit.
Conversely $Z(n)=\Block(Z)(b_n)$, proving degree equality.

Given a total numerical description $D$, first binaryize it. Compute the
majority bit $m(n)$ of $\beta(D)$ on the finite interval $J_n$, with ties
resolved as zero. This is a total $D$-computable sequence. If
$m(n)\ne Z(n)$, at least $|J_n|/2$ positions of that block are errors.
At the end of the block this gives
\begin{equation}\label{eq:blockbound}
 \rho_{b_{n+1}}(\beta(D)\triangle\Block(Z))
 \ge\frac{b_{n+1}-b_n}{2b_{n+1}}\ge\frac14,
\end{equation}
since $b_n/b_{n+1}\le1/2$.
Infinitely many incorrect majority bits would contradict density-zero
error. Thus $m$ and $Z$ differ at only finitely many inputs. Hardcoding
these finitely many corrections yields a $D$-oracle program for $Z$.
This proves $Z\leT D$ without claiming that the finite correction table
can be recovered uniformly from $D$.
\end{proof}

\subsection{Attaching the common information}
Let
\begin{equation}\label{eq:AZ}
 A_Z=\Block(Z)\oplus B_Z,
 \qquad A_Z(2n)=\Block(Z)(n),\quad A_Z(2n+1)=B_Z(n).
\end{equation}
This set is c.e. in $Z$, because its even column is $Z$-computable and its
odd column is c.e. in $Z$.

\begin{lemma}[What every representative computes]\label{lem:anchor}
Every coarse description of $A_Z$ computes $Z$, and no coarse description
of $A_Z$ is $Z$-computable. For $r\in\{\nc,\uc\}$ every
$F\equiv_r A_Z$ satisfies $Z\leT F$, and no such $F$ satisfies $F\leT Z$.
\end{lemma}
\begin{proof}
If $D\approx A_Z$, its even-coordinate projection is a coarse description
of $\Block(Z)$: the number of errors on its first $N$ coordinates is at
most the number of errors of $D$ below $2N$, so its error density is at
most twice that of $D$ there. Lemma~\ref{lem:block} gives $Z\leT D$.
The analogous estimate on odd coordinates makes the odd projection a
coarse description of $B_Z$. A $Z$-computable $D$ would therefore contradict
the relative diagonalization.

Now suppose $F\equiv_r A_Z$. Apply $A_Z\le_{\nc}F$ to the total function
$F$ itself, a description of itself. It computes some $D\in\CD(A_Z)$.
The preceding paragraph gives $Z\leT D\leT F$. If $F\leT Z$, this same
$D$ would be $Z$-computable, a contradiction. The proof allows arbitrary
numerical $F$.
\end{proof}

For the relative fusion theorem use
\begin{equation}\label{eq:anchoredprotect}
 \widetilde P_k=\{2n:n\in\N\}\cup\{2n+1:n\in P_k\},
 \qquad
 \widetilde Q_k=\{2n+1:n\in Q_k\}.
\end{equation}
All of the even column is now protected. On the odd column only finitely
many original columns are protected. These sets are nested and computable;
$\widetilde Q_k$ is infinite with density $2^{-(k+1)}$.
The restriction $A_Z\cap\widetilde P_k$ is $Z$-computable: it consists of
the entire block code and a finite union of the relative computable columns.
Its indices can be selected using $Z''$.

\begin{proof}[Proof of Theorem~\ref{thm:infimum}]
Apply Theorem~\ref{thm:fusion} relative to $Z$, with target $A_Z$ and
protected regions $\widetilde P_k$. It gives a perfect family
$\mathcal P_Z\subseteq\CD(A_Z)$. For distinct $X,Y$ in this family,
any total $F\leT X,Y$ is also below $Z\oplus X,Z\oplus Y$ and hence
below $Z$. On the other hand Lemma~\ref{lem:anchor} gives $Z\leT X,Y$.
Therefore the meet of the degrees of $X,Y$ exists and equals $\degT(Z)$.
They are both strictly above that degree, and in particular are incomparable.
The effective clause and two computable branch labels give such $X,Y\leT Z''$.

Fix $r\in\{\nc,\uc\}$. Lemma~\ref{lem:anchor} says that $\degT(Z)$
is a lower bound for the whole spectrum $\Spec_r(A_Z)$, but is not itself
in that spectrum. Every lower bound for the spectrum is a lower bound
for the chosen pair $X,Y$, which are representatives by
Lemma~\ref{lem:descriptions}. It must therefore be at most their meet
$\degT(Z)$. This proves the greatest-lower-bound assertion and its
nonattainment.
\end{proof}

\section{Which nonuniform classes do have a least representative?}
\label{sec:classification}
The counterexample also suggests an exact positive characterization.
The issue is whether the common information in all descriptions is
represented by a member of the same coarse class.

\begin{definition}
For $f\in\Baire$, its \emph{robust Turing core} is the collection of
degrees
\[
 \Core(f)=\{\mathbf b:(\forall D\in\CD(f))\;\mathbf b\le\degT(D)\}.
\]
This is the degree version of the $X^{\mathfrak c}$ notation in~\cite{HJKS}.
It is a downward-closed collection, not necessarily a single degree.
\end{definition}

\begin{proposition}[Core equals the common lower bounds of representatives]
\label{prop:core}
For $r\in\{\nc,\uc\}$,
\[
 \Core(f)=\{\mathbf b:(\forall\mathbf d\in\Spec_r(f))\;\mathbf b\le\mathbf d\}.
\]
In particular this core is invariant under nonuniform coarse equivalence.
\end{proposition}
\begin{proof}
Every description is a representative, so a common lower bound of all
representatives is below every description. Conversely let
$\mathbf b\in\Core(f)$ and $g\equiv_r f$. The reduction $f\le_{\nc}g$
applied to $g$ itself yields a description $D$ of $f$ with $D\leT g$.
Then $\mathbf b\le\degT(D)\le\degT(g)$. Equality of the collections
follows. Invariance follows by taking $r=\nc$.
\end{proof}

\begin{proposition}[Upward closure of representative spectra]
\label{prop:upward}
For both notions of coarse equivalence, $\Spec_r(f)$ is upward closed in
the Turing degrees.
\end{proposition}
\begin{proof}
Suppose $g\equiv_r f$ and $B\subseteq\N$ satisfies $g\leT B$. Define
\[
 h(2^n)=B(n),\qquad h(m)=g(m)\ \text{when }m\text{ is not a power of }2.
\]
The powers of two have density zero, so $h\approx g$ and
$h\equiv_{\uc}g$. Hence $h\equiv_r f$. Reading its values at powers of
two computes $B$, whereas $B$ computes $h$ because it computes $g$.
Thus $\degT(h)=\degT(B)$, as required. This argument also works for
numerical, rather than binary, $g$.
\end{proof}

\begin{theorem}[Characterization of least nonuniform representatives]
\label{thm:characterization}
For a nonuniform coarse class $\mathcal C$, the following are equivalent:
\begin{enumerate}[label=\textup{(\arabic*)}]
\item $\mathcal C$ has a least Turing representative;
\item its Turing spectrum is a principal upper cone;
\item some $g\in\mathcal C$ satisfies
$g\leT D$ for every $D\in\CD(g)$;
\item $\mathcal C=[\Block(B)]_{\nc}$ for some binary real $B$.
\end{enumerate}
When (4) holds, the least representative degree is $\degT(B)$.
\end{theorem}
\begin{proof}
(1) and (2) are equivalent by Proposition~\ref{prop:upward}; a principal
upper cone includes its base degree. If $g$ is a least representative,
every description of $g$ is another representative, proving (3).

Suppose (3) holds. If $h\in\mathcal C$, the reduction $g\le_{\nc}h$
applied to $h$ itself gives a description $D$ of $g$ with $D\leT h$.
Then $g\leT D\leT h$, so $g$ is least. This already proves (3)$\Rightarrow$(1).

To prove (3)$\Rightarrow$(4), choose a binary graph code $B$ for $g$.
The totality of $g$ gives $B\eqT g$: a value oracle decides graph
membership, and a graph oracle finds the unique value by search. Every
description of $g$ computes $g$, and hence computes $\Block(B)$ itself.
This proves $\Block(B)\le_{\nc}g$. Conversely every description of
$\Block(B)$ computes $B$ by Lemma~\ref{lem:block}, hence computes $g$
itself. Therefore $g\le_{\nc}\Block(B)$, establishing equality of classes.

Finally, if (4) holds, any $h\in\mathcal C$ computes a coarse description
of $\Block(B)$ by applying $\Block(B)\le_{\nc}h$ to $h$ itself.
Such a description computes $B$. Thus $B\leT h$ for every representative,
and $\Block(B)\eqT B$ is itself a representative. It is least, proving (1).
\end{proof}

The map $\mathbf b\mapsto[\Block(B)]_{\nc}$ is indeed an order embedding.
If $B\leT C$, every description of $\Block(C)$ computes $C$, hence
$\Block(B)$; this gives the forward order implication. Conversely, apply
$\Block(B)\le_{\nc}\Block(C)$ to $\Block(C)$ itself and decode the
resulting description of $\Block(B)$ to get $B\leT C$.
The theorem therefore identifies the least-representative classes with
the image of a robust-coding embedding. The argument is included to make
the implication explicit, not to claim that the coarsening-embedding
framework is new; that framework is developed in~\cite{HJKS}.

For the classes constructed here, Proposition~\ref{prop:core} gives the
particularly simple description
\[
 \Core(A_Z)=\{\mathbf b:\mathbf b\le\degT(Z)\},
 \qquad \degT(Z)\notin\Spec_r(A_Z).
\]
Thus even a principal robust core need not have its top degree represented
inside the original coarse class.

\begin{remark}[Uniform classes]
The counterexamples and Proposition~\ref{prop:core} apply to both reduction
notions. Theorem~\ref{thm:characterization} is deliberately stated for
\emph{nonuniform} coarse classes: its block-decoding step hardcodes a finite
correction table depending on the description. No characterization of
uniform classes by this particular code is inferred from that proof.
\end{remark}

\section{Verification boundaries and computational artifacts}
\label{sec:verification}
\subsection{What the construction establishes}
The example refutes the existence of a \emph{least} representative. It
does not show that the spectrum has no minimal elements. Pairwise
incomparability in a perfect subfamily does not by itself determine
minimality in the full upward-closed spectrum.

All statements about lower bounds concern ordinary Turing reductions on
total set or function oracles. The finite-extension proof never substitutes
partial-function degrees. The bound $X_0,X_1\leT\emptyset''$ is an
existence and construction bound using the second jump; it is not an
ordinary executable procedure for evaluating these infinite sets.

The same argument is not a resolution of the effective-dense part of
Gerdes's question. An effective-dense description must explicitly mark
omitted values and may not give any wrong answers. Density-zero changes
therefore do not automatically preserve its description class. The
identity-functional argument of Lemma~\ref{lem:descriptions}, which is
central here, is a specifically coarse argument.

\subsection{Reproducible finite computations}
The archive contains two standard-library Python programs.
\texttt{checks/enumerate\_witness.py} fixes a concrete enumeration of
partial computable functions and produces finite c.e. approximations to
\eqref{eq:A}. The enumeration includes three simple initial programs and
then all finite register-machine programs. Its list and instruction coding,
register-machine semantics, and bounded-execution convention are documented
in the source. A bounded simulation that has not halted is reported as
\emph{unknown}, never as a proof of divergence.

\texttt{checks/check\_finite.py} tests the exact dyadic counting identity,
finite column stopping behavior, monotonicity of sample c.e. approximations,
a finite analogue of the two-reservoir dichotomy, and the density and block
inequalities. Its deterministic JSON output records the test counts.
These tests can catch transcription and implementation mistakes in the
finite ingredients. They cannot establish oracle noncomputability,
density convergence, perfectness, or the minimal-pair theorem. Those
conclusions rest on the proofs above.

\subsection{Proof and source audits}
The accompanying proof ledger records the nonuniform-index distinction,
persistence of each requirement, the order of limits, and the use of
relative jumps. The source ledger records the inspected versions and the
connection with the older cone-avoidance result. There is no Lean code in
this package, and no result is labelled as kernel-checked. Document
compilation and script execution have a narrower meaning than mathematical
proof verification.

The research outcome is therefore a negative answer to the chosen
least-representative assertion, accompanied by a complete conventional
construction and stronger structural statements. Its relationship to the
existing theorem is explicit; whether the stronger formulations warrant
a claim of new mathematics requires a specialist literature and proof review.

\clearpage
\appendix
\section{Construction pseudocode}\label{app:pseudocode}
This pseudocode describes the level construction, not an ordinary
oracle-free implementation. \texttt{ExistsCrossDisagreement} is decided
using $Z'$ after actual indices for the protected $Z$-computable pieces
have been supplied.
\begin{lstlisting}
level[empty] := empty string
L[0] := 0
for s = 0, 1, 2, ...:
    children := split every node at a fresh position in Q[s]
                using admissible extensions at protection level s
    for each ordered pair of distinct child labels (u, v):
        for e, j <= s:
            if ExistsCrossDisagreement(children[u], children[v],
                                       s, e, j):
                (alpha, beta, n) := first finite witness
                children[u] := alpha
                children[v] := beta
            else:
                record: every common total output for this task
                        is computable in Z
    pad all children admissibly to a common length L[s+1] > L[s]
    level := children
    raise protection to P[s+1] beyond the current prefixes
\end{lstlisting}
For any fixed branch label, the construction never revises one of its
bits. It may use noncomputable decisions to decide which finite extension
to take. These are compatible assertions.

\begin{thebibliography}{9}
\raggedright
\addcontentsline{toc}{section}{References}
\bibitem{Gerdes}
Peter M. Gerdes.
\emph{Comparing Notions of Dense Computability on $\omega^\omega$ and
$2^\omega$}. arXiv:2508.06925v1, 9 August 2025.
Question 7, printed p.~64.
\url{https://arxiv.org/abs/2508.06925}.

\bibitem{GerdesSlides}
Peter M. Gerdes.
\emph{Comparing Notions of Robust Computability on $\omega^\omega$ and
$2^\omega$}. ASL North American Meeting, July 2026.
Slide 45 of 48; final overlay at PDF page 93 of 98.
Author-hosted slides inspected 18 September 2026.
\url{https://invariant.org/coarse.pdf}.

\bibitem{HJKS}
Denis R. Hirschfeldt, Carl G. Jockusch, Jr., Rutger Kuyper, and Paul E. Schupp.
\emph{Coarse Reducibility and Algorithmic Randomness}.
Journal of Symbolic Logic \textbf{81}(3) (2016), 1028--1046.
\href{https://doi.org/10.1017/jsl.2015.70}{doi:10.1017/jsl.2015.70}.
Preprint arXiv:1505.01707v1, 7 May 2015; Definition 3.1,
Theorems 4.2 and 4.3, and the coarsening framework of Section 2.
\url{https://arxiv.org/abs/1505.01707}.

\bibitem{Survey}
\emph{Turing Degrees: A Unified Research Report}.
Prepared for Vladimir Reshetnikov, 18 September 2026.
User-supplied source \texttt{turing\_degrees\_unified.tex};
entry C1 and its minimal-pair diagnostic.
\end{thebibliography}
\end{document}
